\section{Discussion}
\subsection{Authorization location is the main security result}
The most important distinction in these experiments is not between fuzzy and deterministic logic, but between \emph{policy after retrieval} and \emph{policy before retrieval}. The prompt-only system contains an authorization instruction, yet its retriever still has access to the complete corpus. Consequently, every unauthorized request crosses the retrieval boundary even though only five canary strings are reproduced. This gap matters operationally: once a protected record is in context, security relies on the behavior of a probabilistic generator and on the completeness of output monitoring. Pre-retrieval authorization instead makes the candidate set itself a policy object. Under the modeled assumptions, an unauthorized record is absent before model invocation.

This ordering also clarifies the role of risk scoring. A risk controller can reduce workload, trigger step-up authentication, or reject suspicious activity earlier, but it should not decide which protected records a principal is allowed to retrieve. Our held-out experiment makes that separation empirical rather than rhetorical. Lexical suspicion features fail under paraphrase, while ACL denial remains unchanged. A detector may therefore be useful for operational adaptation without being part of the confidentiality proof obligation.

\subsection{Utility depends on retrieval composition, not only policy strictness}
The original $k=2$ result could easily have produced the misleading conclusion that security caused a 15-point utility loss. The paired diagnostic shows a more specific mechanism: the target was present for all legitimate queries, but each architecture paired it with a different distractor. For the small instruction-tuned model used here, the identity of that second record substantially affected answer success. The controlled $k=1$ experiment is therefore more than a robustness check; it demonstrates why security ablations must control retrieval composition before interpreting answer-quality differences causally.

This finding does not imply that narrower authorization scopes always preserve utility. In a real system, a user may legitimately need information spread across several records, and a policy can remove relevant evidence. The correct design principle is instead to compare architectures at matched information availability when estimating the \emph{overhead of enforcement}, and separately study policy-induced information loss when that is the research question.

\subsection{What the 100K study adds}
The scale study contributes a systems perspective rather than another authorization-correctness test. Because the ACL check is set membership, correctness does not become more or less true at 100K patients. What changes is the retrieval environment: candidate pools become larger, latency rises, and the best distractor becomes more similar to the target. The proportional-ACL condition is especially informative because its authorized scope grows with the corpus; it shows that pre-filtering is not a magical constant-time operation when user entitlements themselves scale. Conversely, the fixed-ACL and target-fixed modes are intentionally invariant controls and should remain labeled as such.

The LLM subset gives only limited evidence about generation under scale. Unfiltered answer success declines numerically at 100K but not significantly, and proportional ACL was not passed through the generator. We therefore avoid a claim that larger corpora necessarily reduce generation quality. What is supported is the retrieval-side observation that larger candidate sets increase cost and interference, while precise target scoping provides a high-utility control for the explicitly targeted queries used here.
